\chapter{Entrelacement dodécaphase--Witt et lecture exacte de \texorpdfstring{\(3/11\)}{3/11}}
\label{chap:entrelazador-accion}
\index{opérateur d'entrelacement!dodécaphase--Witt}
\index{Witt!module d'incidence}
\index{survie!projecteur de rang trois}

L'apparition répétée du quotient \(3/11\) dans l'écran dodécaphasique,
dans la polarité d'une étoile de Witt et dans la filtration nonadique pose
une question plus forte que l'égalité arithmétique de trois fractions. La
question correcte est opératorielle : déterminer si ces trois apparitions sont
des traces du même secteur positif vu dans des réalisations différentes et, si elles le
sont, construire les applications qui les entrelacent sans introduire de donnée
métrologique ultérieure.

La réponse comporte quatre niveaux. Premièrement, l'incidence point--hexade
construit une isométrie canonique entre le module dodécaphasique de dimension
onze et un espace propre de dimension onze du plan de Witt. Deuxièmement,
l'étoile de Witt agit sur ce même module, mais comme une involution
signée ; son \(3/11\) est un indice, non la trace d'un projecteur positif.
Troisièmement, le secteur orienté de survie apparaît dans \(t=7\), où un
unique espace de onze branches contient un secteur de rang trois. Quatrièmement, la
droite d'orientation de ce triplet, avec les témoins internes
\(B_K\) et \(u_K\), sélectionne de manière variationnelle un unique normalisateur
\(S_3\), fixe son orientation et produit par moyenne de Reynolds et
décomposition polaire l'entrelacement isométrique secteur--dodécaphase. Le
passage dodécaphase--Witt est une seconde isométrie d'incidence. Dans le calibre
étiqueté initial, les deux flèches ne réalisent pas encore une unique représentation
de \(S_3\) ; la recalibration conjointe construite à la fin du chapitre transporte
la dodécade, l'origine, l'orientation, la chambre et la mémoire, et établit cette
équivariance sans modifier le no-go de l'étiquetage fixe.
Les moyennes de Reynolds, les décompositions en représentations et le
facteur polaire utilisés dans cette section sont les standards
\cite{Serre1977Representations,HornJohnson2013}.

\begin{figure}[htbp]
\centering
\begin{tikzpicture}[
  font=\footnotesize,
  box/.style={draw,rounded corners=2mm,align=center,minimum width=2.35cm,minimum height=.98cm},
  arr/.style={-{Stealth[length=2.2mm]},thick},
  link/.style={<->,thick,dashed}
]
  \node[box,draw=hmtblue,fill=hmtlightblue] (seed) at (0,2.2)
    {germes\\\(104\,976\)};
  \node[box,draw=hmtblue,fill=hmtlightblue] (u) at (3,2.2)
    {émissions\\\(\mathcal U_6,\ |\mathcal U_6|=468\)};
  \node[box,draw=hmtgold,fill=hmtlightgold] (ppi) at (6,2.2)
    {microfibra\\\(\Pi_\pi^{(468)},\ \operatorname{rank}=5\)};
  \node[box,draw=hmtred,fill=hmtlightred] (d) at (9,2.2)
    {défaut positif\\\(D_{\rm switch}\)};
  \node[box,draw=hmtviolet,fill=hmtlightviolet] (lam) at (12,2.2)
    {capacidad residual\\\(\lambda_*=1-\delta_*\)};

  \draw[arr,hmtblue] (seed)--node[above=16pt,font=\scriptsize,fill=white,inner sep=1pt]{\(F\)} (u);
  \draw[arr,hmtgold] (u)--node[above=16pt,font=\scriptsize,fill=white,inner sep=1pt]{\(\tau_{468}\)} (ppi);
  \draw[arr,hmtred] (ppi)--node[above=16pt,font=\scriptsize,fill=white,inner sep=1pt]{\(D_{\rm switch}\)} (d);
  \draw[arr,hmtviolet] (d)--node[above=16pt,font=\scriptsize,fill=white,inner sep=1pt]{\(\tau+\mathrm{CR}\)} (lam);

  \node[box,draw=hmtgold,fill=hmtlightgold] (g) at (1.5,-.2)
    {secteur orienté \(t=7\)\\\(L_{\rm or}\otimes(H_G,P_G)\)};
  \node[box,draw=hmtgreen,fill=hmtlightgreen] (v) at (5,-.2)
    {dodécaphase\\\(V_{11},P_3\)};
  \node[box,draw=hmtgreen,fill=hmtlightgreen] (e) at (8.5,-.2)
    {incidence de Witt\\\(E_{30},Q_3\)};
  \node[box,draw=hmtgray,fill=hmtpaper] (t) at (12,-.2)
    {lecteur spectral\\\(\mathcal L_{\beta,*}\to T_{\rm band}\)};

  \draw[arr,hmtgold] (g)--node[above=16pt,font=\scriptsize,fill=white,inner sep=1pt]{\(W_{GK}\)} (v);
  \draw[arr,hmtgreen] (v)--node[above=16pt,font=\scriptsize,fill=white,inner sep=1pt]{\(U_W=I/6\)} (e);
  \draw[arr,hmtgray] (lam)--(t);
  \draw[arr,hmtred] (v) to[bend left=20]
    node[below=3pt,font=\scriptsize,fill=white,inner sep=1pt]{\((3/11)\Delta_4\)} (d);
\end{tikzpicture}

\caption{Architecture de la fermeture prédimensionnelle. La ligne supérieure construit
la mesure et l'action à partir du quotient des émissions ; la ligne inférieure transporte
le secteur orienté \(t=7\) vers la dodécaphase, puis la dodécaphase vers
Witt. Les flèches continues sont des isométries de projecteurs démontrées. La
première est fixée par la sélection variationnelle du
\cref{thm:identificacion-isometrica-tres-proyectores} ; la seconde n'est pas
\(H_*\)-équivariante dans le calibre fixe et le devient au moyen du
changement intégral de calibre du \cref{thm:recalibracion-hstar-witt}.}
\label{fig:entrelazador-accion}
\end{figure}

\section{Typage des objets associés à \texorpdfstring{\(3/11\)}{3/11}}

Soit
\[
 V_{12}=\mathbb R^{12},
 \qquad
 V_{11}=\mathbf 1^\perp
 =\left\{v\in\mathbb R^{12}:\sum_{p=1}^{12}v_p=0\right\}.
\]
L'action icosaédrique déclarée sur les douze positions se décompose comme
\[
 V_{12}=V_1\oplus V_3\oplus V_{3'}\oplus V_5.
\]
Notons \(P_3\) le projecteur orthogonal sur le premier terme
tridimensionnel. Alors
\[
 P_3^2=P_3=P_3^*,
 \qquad
 \operatorname{rank}P_3=3,
 \qquad
 \tau_{11}(P_3)
 =\frac1{11}\operatorname{Tr}_{V_{11}}P_3
 =\frac3{11}.
\]

Il y a trois objets proches et un quatrième objet de type distinct :
\begin{enumerate}
  \item l'incidence point--hexade transporte \(P_3\) vers un projecteur
  \(Q_3\) de rang trois sur un module de Witt ;
  \item la survie dans \(t=7\) distingue trois branches parmi onze et
  produit un projecteur \(P_G\) de trace normalisée \(3/11\) ;
  \item un opérateur d'entrelacement orienté identifie
  \(P_3V_{11}\) à la réalisation orientée de \(P_GH_G\) ;
  \item une étoile de Witt a sept modes pairs et quatre impairs,
  de sorte que sa trace normalisée vaut également \(3/11\), mais elle contient un spectre
  négatif et n'est pas un projecteur.
\end{enumerate}
La distinction entre projecteur et involution sera essentielle pour que la loi
d'action de la section suivante soit positive.

\section{La matrice d'incidence point--hexade}

Soit \(\mathcal H\) l'ensemble des \(132\) hexades du plan
\[
 W_{12}=S(5,6,12).
\]
On définit
\[
 I:\mathbb R^{12}\longrightarrow\mathbb R^{\mathcal H},
 \qquad
 (Iv)(H)=\sum_{p\in H}v_p,
\]
ou, de manière équivalente,
\[
 I_{H,p}=\mathbf 1_{\{p\in H\}}.
\]
Cet opérateur n'enregistre que l'incidence. Il ne contient ni \(\alpha\), ni unités,
ni masses, ni CKM, ni Barbero, ni aucun objectif physique.

\begin{lemma}[Identité métrique de l'incidence]
\label{lem:incidencia-metrica-witt}
On a
\[
 \boxed{I^*I=36I_{12}+30J_{12}.}
\]
Par conséquent, pour tout \(v\in V_{11}\),
\[
 \|Iv\|^2=36\|v\|^2.
\]
\end{lemma}

\begin{proof}
Dans un système \(S(5,6,12)\), chaque point appartient à
\[
 r=\frac{\binom{11}{4}}{\binom54}=66
\]
hexades, et chaque paire de points appartient à
\[
 \lambda_2=\frac{\binom{10}{3}}{\binom43}=30
\]
hexades. Les entrées diagonales de \(I^*I\) valent \(66\) et les entrées non
diagonales valent \(30\). Donc
\[
 I^*I=(66-30)I_{12}+30J_{12}.
\]
Comme \(J_{12}\) s'annule sur \(\mathbf1^\perp\), on obtient la seconde
identité.
\end{proof}

Pour identifier la branche spectrale sélectionnée par \(I\), nous introduisons
\[
 G_{H,H'}=\binom{|H\cap H'|}{4}.
\]

\begin{lemma}[Identité spectrale de l'incidence]
\label{lem:incidencia-espectral-witt}
Avec \(J_{132,12}\) la matrice rectangulaire de uns,
\[
 \boxed{GI=30I+15J_{132,12}.}
\]
En particulier, si \(v\in V_{11}\), alors
\[
 G(Iv)=30Iv.
\]
\end{lemma}

\begin{proof}
L'entrée \((H,p)\) de \(GI\) est
\[
 \sum_{H'\ni p}\binom{|H\cap H'|}{4}.
\]
Le recensement exhaustif de \(W_{12}\) donne \(45\) lorsque \(p\in H\) et \(15\)
lorsque \(p\notin H\). Par conséquent, l'entrée vaut
\(30I_{H,p}+15\). La reconstruction des \(132\) hexades à partir du code
de Golay ternaire étendu vérifie les \(132\cdot12\) entrées comme des
entiers.
\end{proof}

Définissons
\[
 E_{30}^{\rm inc}:=I(V_{11})\subseteq\ker(G-30I).
\]
Sa dimension est onze par le \cref{lem:incidencia-metrica-witt}. Le calcul
modulo \(1009\) obtient
\[
 \operatorname{rank}_{\mathbb F_{1009}}(G-30I)=121.
\]
Comme la réduction modulo \(1009\) donne une borne inférieure du rang
rationnel, on en déduit
\(\dim_{\mathbb Q}\ker(G-30I)\le11\). L'inclusion précédente fournit
onze vecteurs linéairement indépendants dans ce noyau. Par conséquent,
\[
 E_{30}^{\rm inc}=\ker(G-30I).
\]
Nous écrirons désormais \(E_{30}\).

\begin{theorem}[Opérateur d'entrelacement canonique dodécaphase--Witt]
\label{thm:entrelazador-dodecafase-witt}
L'opérateur
\[
 \boxed{
 U_W=\frac16 I\big|_{V_{11}}:V_{11}\longrightarrow E_{30}
 }
\]
est une isométrie surjective. Pour tout
\(g\in\operatorname{Aut}(W_{12})\cong M_{12}\),
\[
 U_W\rho_{12}(g)=\rho_{\mathcal H}(g)U_W.
\]
Parmi les noyaux point--hexade équivariants, la condition d'isométrie et
le signe positif sur les incidences déterminent \(U_W\) de manière unique.
\end{theorem}

\begin{proof}
L'isométrie découle de \(I^*I=36I\) sur \(V_{11}\). Son image est de
dimension onze et est contenue dans \(E_{30}\), également de dimension onze ;
elle est donc surjective. L'équivariance traduit
\[
 p\in H\quad\Longleftrightarrow\quad g(p)\in g(H).
\]

Pour l'unicité, \(M_{12}\) est transitif sur les drapeaux
\((p,H)\), avec \(p\in H\), et sur les antidrapeaux. Tout noyau
équivariant a deux valeurs, \(a\) sur les drapeaux et \(b\) sur les
antidrapeaux :
\[
 T=aI+b(J-I)=(a-b)I+bJ.
\]
Le terme \(J\) s'annule sur \(V_{11}\), de sorte que \(T\) y est un
multiple de \(I\). L'isométrie impose \(|a-b|=1/6\) et la normalisation
positive sélectionne \(a-b=1/6\).
\end{proof}

\section{Transport positif du projecteur tritique}

L'opérateur d'entrelacement permet de définir
\[
 \boxed{Q_3=U_WP_3U_W^*.}
\]
On obtient
\[
 Q_3^2=Q_3=Q_3^*,
 \qquad
 \operatorname{rank}Q_3=3,
\]
et le diagramme
\[
\begin{CD}
V_{11} @>{U_W}>> E_{30}\\
@V{P_3}VV @VV{Q_3}V\\
V_{11} @>{U_W}>> E_{30}
\end{CD}
\]
commute. En particulier,
\[
 \boxed{\tau_{V_{11}}(P_3)=\tau_{E_{30}}(Q_3)=\frac3{11}.}
\]

La canonicité de \(U_W\) appartient au plan étiqueté et à son action
d'automorphismes. Il faut la distinguer de l'écran \(A_5/C_5\) utilisé pour
définir \(P_3\) : dans la réalisation étiquetée construite, seule l'identité
de cette action concrète de \(A_5\) préserve l'ensemble fixe des hexades. Ainsi,
\(U_W\) transporte canoniquement le \(P_3\) déjà donné vers
\(Q_3\), mais ne transforme pas à lui seul l'action \(A_5/C_5\) en une
sous-action de \(M_{12}\). Le pont avec le secteur orienté sera construit séparément
à partir des données marquées \(K\) et \(B_K\).

\section{Indice signé de l'étoile de Witt}

Soit \(B\) une tétrade. Les quatre hexades qui la contiennent s'écrivent
\[
 H_i=B\sqcup P_i,\qquad |P_i|=2,
\]
où les quatre pétales \(P_i\) partitionnent le complémentaire de \(B\).
L'involution d'étoile \(R_B\) fixe \(B\) et échange les deux points
de chaque pétale. Son type de cycles sur douze points est \(1^4\,2^4\). Dans
\(V_{11}\), et au moyen de \(U_W\) également dans \(E_{30}\), son spectre est
\[
 \operatorname{Spec}(R_B)=\{1^{[7]},(-1)^{[4]}\}.
\]
Par conséquent,
\[
 \boxed{\tau_{11}(R_B)=\frac{7-4}{11}=\frac3{11}.}
\]

\begin{theorem}[Obstruction projecteur--étoile]
\label{thm:no-go-proyector-estrella}
Il n'existe pas d'isomorphisme \(F:V_{11}\to E_{30}\) tel que
\[
 FP_3=R_BF.
\]
\end{theorem}

\begin{proof}
S'il existait, \(P_3\) et \(R_B\) seraient semblables. Cependant,
\[
 \operatorname{Spec}(P_3)=\{1^{[3]},0^{[8]}\},
 \qquad
 \operatorname{Spec}(R_B)=\{1^{[7]},(-1)^{[4]}\}.
\]
Leurs polynômes minimaux sont \(x(x-1)\) et
\((x-1)(x+1)\), respectivement. La similitude est impossible.
\end{proof}

L'égalité \(3/11\) possède ainsi deux lectures sur le même module. Pour
\(Q_3\), c'est la masse d'un secteur positif de rang trois. Pour \(R_B\), c'est
un indice de parité. Dans une action positive, \(Q_3\) doit intervenir ;
l'étoile conserve l'orientation et distingue les feuillets, mais ne remplace pas le
projecteur.

\section{Réduction homogène \texorpdfstring{\(11\to3\to1\)}{11-3-1} en \texorpdfstring{\(t=7\)}{t=7}}

La lecture précédente de \(3/11\) dans \(t=8\) comparait trois racines locales
à onze continuations d'une seule racine. Le numérateur et le dénominateur
n'appartenaient pas au même espace. Le tableau détaillé situe l'objet correct
en \(t=7\) : il y a onze branches locales ; les onze survivent à un horizon, trois
survivent à deux et une survit à trois,
\[
 \boxed{11\longrightarrow3\longrightarrow1.}
\]

\begin{center}
\small
\begin{tabularx}{\textwidth}{cYcc}
\toprule
indice&triplet de mots&\(h=2\)&\(h=3\)\\
\midrule
0&\texttt{200112|212221|110110}&3&0\\
1&\texttt{200121|212212|110110}&3&1\\
2&\texttt{200211|212122|110110}&3&0\\
3&\texttt{202122|210211|110110}&0&0\\
4&\texttt{202212|210121|110110}&0&0\\
5&\texttt{202221|210112|110110}&0&0\\
6&\texttt{210111|202222|110110}&0&0\\
7&\texttt{210222|202111|110110}&0&0\\
8&\texttt{212112|200221|110110}&0&0\\
9&\texttt{212121|200212|110110}&0&0\\
10&\texttt{212211|200122|110110}&0&0\\
\bottomrule
\end{tabularx}
\end{center}

Soit \(\mathcal B_7\) l'ensemble de ces onze branches et
\[
 H_G=\ell^2(\mathcal B_7).
\]
On définit
\[
 P_G\delta_b=
 \begin{cases}
 \delta_b,&b\text{ admet un prolongement jusqu’à l’horizon deux},\\
 0,&\text{sinon}.
 \end{cases}
\]
Alors
\[
 P_G^2=P_G=P_G^*,
 \qquad
 \operatorname{rank}P_G=3,
 \qquad
 \boxed{\tau_{H_G}(P_G)=\frac3{11}.}
\]
C'est la troisième réalisation positive.

\section{L'action résiduelle de \texorpdfstring{\(S_3\)}{S3} et les régions positives de \texorpdfstring{\(\pi\)}{pi}}

Une permutation des trois positions finales, appliquée simultanément aux
deux premiers mots de chaque branche en laissant le troisième fixe, préserve
\(\mathcal B_7\). Cette règle définit une action de \(S_3\). Ses orbites
ont pour tailles
\[
 3+3+1+1+3=11.
\]
Le support de \(P_G\), formé des indices \(0,1,2\), est une orbite
complète, de sorte que \(P_G\) commute avec \(S_3\).

Le caractère de la représentation sur les onze branches, lu sur l'identité,
la transposition et le cycle d'ordre trois, est
\[
 \chi_{\mathcal B_7}=(11,5,2)
 =5\chi_{\rm triv}+3\chi_{\rm std}.
\]
Sur l'image de \(P_G\),
\[
 \chi_G=(3,1,0)
 =\chi_{\rm triv}+\chi_{\rm std}.
\]

Les trois régions positives de la microfibre de \(\pi\) ont pour signatures
\[
 339555,\qquad393555,\qquad933555.
\]
Les queues du premier mot des trois branches survivantes sont
\[
 112,\qquad121,\qquad211.
\]
La substitution en coordonnées \(1\mapsto3\), \(2\mapsto9\) donne
\[
\begin{array}{c|c|c}
\text{branche}&\text{queue}&\text{région positive}\\ \hline
0&112&339555\\
1&121&393555\\
2&211&933555.
\end{array}
\]

\begin{theorem}[Pont horizontal--vertical de la microfibre positive]
\label{thm:puente-pi-puerta-s3}
La règle précédente est une bijection \(S_3\)-équivariante entre les trois
régions positives de \(\pi\) et les trois branches sélectionnées par \(P_G\).
L'application conserve la position de l'unique symbole exceptionnel.
\end{theorem}

\begin{proof}
Les deux ensembles sont les trois mots obtenus en plaçant un unique symbole
exceptionnel dans l'une des trois positions. La substitution s'applique coordonnée par
coordonnée et commute donc avec toute permutation des positions. La
position exceptionnelle retrouve de manière univoque l'élément d'origine.
\end{proof}

\section{Dépendance à l'orientation dans l'opérateur d'entrelacement dodécaphasique--Witt}

Pour tout normalisateur \(H=N_{A_5}(C_3)\simeq S_3\) d'un sous-groupe
d'ordre trois dans \(A_5\), nous définissons le caractère abstrait du quotient
\[
 \operatorname{sgn}_H:H\longrightarrow H/C_3\simeq C_2
 \longrightarrow\{+1,-1\},
 \qquad
 \operatorname{sgn}_H(h)=
 \begin{cases}
 +1,&h\in C_3,\\
 -1,&h\in H\setminus C_3.
 \end{cases}
\]
Ce caractère n'est pas le signe ambiant d'une permutation de cinq lettres :
tout élément de \(A_5\) a pour signe ambiant \(+1\). C'est le caractère
intrinsèque de la réflexion dans le quotient \(H/C_3\). Avec cette convention, le
caractère du secteur icosaédrique restreint à \(H\) est
\[
 \chi_{P_3V_{11}|H}=(3,-1,0)
 =\chi_{\operatorname{sgn}_H}+\chi_{\rm std}.
\]
En revanche,
\[
 \chi_{\operatorname{im}P_G}=(3,1,0)
 =\chi_{\rm triv}+\chi_{\rm std}.
\]

\begin{corollary}[Obstruction à l'opérateur d'entrelacement non orienté]
\label{cor:obstruccion-entrelazador-incidencia-p3}
Il n'existe pas d'isomorphisme \(S_3\)-équivariant
\[
 \operatorname{im}P_G\longrightarrow P_3V_{11}
\]
pour les actions non tordues.
\end{corollary}

\begin{proof}
Des représentations isomorphes possèdent le même caractère. Les valeurs sur une
transposition sont \(1\) et \(-1\).
\end{proof}

Soit \((e_0,e_1,e_2)\) la base positionnelle de
\(\mathbb R^3_{\rm perm}\). Nous définissons
\[
 L_{\rm or}
 =\bigwedge\nolimits^3\mathbb R^3_{\rm perm},
 \qquad
 \ell_{\rm or}=e_0\wedge e_1\wedge e_2
\]
la droite d'orientation du triplet, munie de l'ancrage positif
\(\ell_{\rm or}\). L'action de \(S_3\) sur cette droite est précisément son
caractère de signe abstrait. En tensorisant par elle,
\[
 \chi_{L_{\rm or}\otimes\operatorname{im}P_G}
 =(3,-1,0),
\]
de sorte que l'obstruction des caractères disparaît. L'égalité des
caractères démontre l'existence d'une classe d'isomorphismes, mais n'en
sélectionne pas encore un. La sélection s'obtient à partir de deux témoins déjà présents dans la
dodécaphase.

\section{Sélection variationnelle du normalisateur \texorpdfstring{\(S_3\)}{S3}}

Soit
\[
 K=(234,543,140,729,659,824,621,58,914,794,146,601)
\]
la carte décimale du sceau. Le témoin commun minimal des deux lectures est
\[
 B_K=H_e\cap P_\pi
 =\{p:K_p\ge729\}
 =\{4,6,9,10\},
\]
dans la numérotation mathématique \(1,\ldots,12\). Dans la convention équivalente
avec des indices commençant à zéro, le même ensemble est \(\{3,5,8,9\}\). Nous définissons deux
vecteurs du secteur \(P_3\) :
\[
 b_K=P_3\left(\mathbf1_{B_K}-\frac{|B_K|}{12}\mathbf1\right),
 \qquad
 u_K=P_3\left(K-\overline K\,\mathbf1\right).
\]
Le premier satisfait exactement
\[
 \|b_K\|^2=1.
\]

Le groupe \(A_5\) contient dix sous-groupes \(C_3\). Pour chacun, soit
\[
 H(C_3)=N_{A_5}(C_3)\simeq S_3
\]
et soit le projecteur de signe
\[
 e_{{\rm sgn},H}
 =\frac16\sum_{h\in H}\operatorname{sgn}_H(h)\rho(h).
\]
Ici \(\operatorname{sgn}_H\) est le caractère abstrait de \(H/C_3\)
défini ci-dessus, non le signe ambiant —trivial— de \(A_5\).
L'énergie d'orientation associée au témoin commun est
\[
 \mathcal E_B(H)=\|e_{{\rm sgn},H}b_K\|^2.
\]
Le recensement exact dans \(\mathbb Q(\sqrt5)\) produit un maximum unique :
\[
 \mathcal E_B(H_*)
 =\frac23+\frac{2}{15}\sqrt5.
\]
Le deuxième niveau est
\[
 \frac13+\frac{2}{15}\sqrt5,
\]
de sorte que l'écart exact est
\[
 \boxed{\mathcal E_B(H_*)-\mathcal E_B(H_{\rm segundo})=\frac13.}
\]
Le même \(H_*\) est le maximum unique lorsque \(b_K\) est remplacé par \(u_K\).
Le sélecteur ne dépend donc pas d'un seul codage du témoin.

Dans cet étiquetage,
\[
 H_*=N_{A_5}\bigl(\langle c_*\rangle\bigr),
\qquad
 c_*=(2\,3\,4).
\]
La règle est naturelle sous un réétiquetage simultané par \(A_5\) : en
transportant \(B_K\), \(K\) et \(P_3\), le maximum est conjugué par le même
élément.

\section{Sélection de l'orientation et présentation de \texorpdfstring{\(S_3\)}{S3}}

Dans \(H_*\), considérons
\[
 s(g)=\langle b_K,\rho(g)u_K\rangle.
\]
Le maximum parmi les deux éléments d'ordre trois sélectionne de manière univoque
\[
 c_*=(0,1,3,4,2),
\]
dans la notation des images de la permutation de cinq points, et le maximum parmi
les trois involutions sélectionne
\[
 r_*=(1,0,4,3,2).
\]
On vérifie exactement
\[
 r_*c_*r_*=c_*^{-1}.
\]
La présentation du secteur orienté au moyen du cycle de positions
\(c=(0\,1\,2)\) et de la réflexion \(r=(0\,2)\), qui fixe la position centrale,
détermine alors un unique isomorphisme
\[
 \theta:S_3^{\rm or}\xrightarrow{\sim}H_*,
 \qquad
 \theta(c)=c_*,
 \qquad
 \theta(r)=r_*.
\]
Ici \(S_3^{\rm or}\) désigne le groupe des réétiquetages du secteur orienté
\(t=7\), engendré par \(c\) et \(r\) ; il n'introduit pas de classe opératoire
supplémentaire.
L'unique branche qui atteint l'horizon trois est précisément la branche centrale
\(1\) ; cette donnée fixe laquelle des trois positions reste immobile sous la
réflexion.

\section{Moyenne de Reynolds et facteur polaire}

Soit
\[
 W_{\rm or}=L_{\rm or}\otimes\mathbb R^3_{\rm perm}
\]
le triplet du secteur orienté. L'ancrage étant déjà fixé, nous prenons
\[
 \boxed{w=\ell_{\rm or}\otimes e_1,}
\]
où \(e_1\) représente la position centrale, la seule qui atteint l'horizon
trois. Nous définissons
\[
 A
 =\sum_{g\in S_3}
 \rho_{W_{\rm or}}(g)\,
 |w\rangle\langle b_K|\,
 \rho_{V_{11}}(\theta(g))^{-1}.
\]
Par construction,
\[
 A\rho_{V_{11}}(\theta(g))
 =\rho_{W_{\rm or}}(g)A,
\qquad
 A=AP_3.
\]
Le calcul exact dans \(\mathbb Q(\sqrt5)\) obtient
\[
 \operatorname{rank}A=3.
\]
De plus, \(G_A=AA^*\) a toutes ses entrées diagonales égales à
\[
 3+\frac25\sqrt5
\]
et toutes ses entrées non diagonales égales à
\[
 \frac52+\frac35\sqrt5.
\]
Ses deux valeurs propres isotypiques sont
\[
 \lambda_{\rm sgn}=8+\frac85\sqrt5,
 \qquad
 \lambda_{\rm std}=\frac12-\frac15\sqrt5.
\]
Toutes deux sont strictement positives. En particulier, \(G_A^{-1/2}\) existe et
la décomposition polaire
\[
 \boxed{U_{GK}=G_A^{-1/2}A}
\]
est une isométrie équivariante de \(P_3V_{11}\) sur \(W_{\rm or}\) :
\[
 U_{GK}^*U_{GK}=P_3,
 \qquad
 U_{GK}U_{GK}^*=I_{W_{\rm or}}.
\]

Soit
\[
 J_G:W_{\rm or}\longrightarrow L_{\rm or}\otimes H_G
\]
l'inclusion isométrique dans les trois branches sélectionnées. Alors
\[
 J_GJ_G^*=I_{L_{\rm or}}\otimes P_G.
\]
L'isométrie partielle secteur--dodécaphase est notée
\[
 \boxed{
 W_{GK}=U_{GK}^*J_G^*
 :L_{\rm or}\otimes H_G\longrightarrow V_{11}
 }
\]
satisfait
\[
 W_{GK}^*W_{GK}
 =I_{L_{\rm or}}\otimes P_G,
\qquad
 W_{GK}W_{GK}^*=P_3,
\]
et, par conséquent,
\[
 P_3W_{GK}
 =W_{GK}
 =W_{GK}(I_{L_{\rm or}}\otimes P_G).
\]

\begin{theorem}[Identification isométrique relative de trois projecteurs]
\label{thm:identificacion-isometrica-tres-proyectores}
Relativement à l'action \(A_5/C_5\), au projecteur \(P_3\), au sceau
\(K\), au témoin doublement construit \(B_K\), au secteur fini orienté
\(t=7\), au caractère abstrait \(\operatorname{sgn}_{H_*}\), à l'ancrage
\(\ell_{\rm or}\) et à la règle variationnelle \(\mathcal E_B\), il existe une
unique isométrie partielle normalisée \(W_{GK}\) qui identifie \(P_G\) à
\(P_3\). Indépendamment, l'isométrie d'incidence \(U_W\) identifie
\(P_3\) à \(Q_3\). Par conséquent, on obtient la chaîne suivante
d'identifications isométriques de projecteurs :
\[
\begin{CD}
L_{\rm or}\otimes H_G
 @>{W_{GK}}>>
 V_{11}
 @>{U_W}>>
 E_{30}\\
@V{I_{L_{\rm or}}\otimes P_G}VV
 @VV{P_3}V
 @VV{Q_3}V\\
L_{\rm or}\otimes H_G
 @>{W_{GK}}>>
 V_{11}
 @>{U_W}>>
 E_{30}.
\end{CD}
\]
Les traces typées satisfont
\[
 \tau_{L_{\rm or}\otimes H_G}(I_{L_{\rm or}}\otimes P_G)
 =\tau_{V_{11}}(P_3)
 =\tau_{E_{30}}(Q_3)
 =\frac3{11}.
\]
L'énoncé, dans le calibre fixe, identifie isométriquement trois
projecteurs ; il n'affirme pas encore que leurs trois espaces soient une même
représentation d'un groupe commun.
\end{theorem}

\begin{proof}
L'énergie \(\mathcal E_B\) sélectionne un unique normalisateur \(H_*\) ;
le critère de score \(s(g)\) sélectionne une unique paire \((c_*,r_*)\), et avec elle un
unique isomorphisme \(\theta\). La moyenne de Reynolds est alors
équivariante pour les actions de \(S_3^{\rm or}\) et \(H_*\) reliées
par \(\theta\), et possède le rang trois. La partie positive de sa décomposition
polaire fixe la normalisation sans introduire de base. Les identités
partielles précédentes démontrent que \(W_{GK}\) identifie
\(P_G\) et \(P_3\). D'autre part, le
\cref{thm:entrelazador-dodecafase-witt} identifie \(P_3\) et \(Q_3\)
comme projecteurs au moyen de \(U_W\). Ces deux identités opératorielles font
commuter ce diagramme de calibre fixe, sans y étendre l'équivariance de la
première flèche à la seconde.
\end{proof}

\begin{proposition}[Obstruction à l'équivariance triple dans le calibre fixé]
\label{prop:no-go-equivarianza-hstar-witt}
La construction précédente ne démontre pas que \(U_W\) soit
\(H_*\)-équivariante. Dans l'étiquetage déclaré, l'action fixe de
\(A_5/C_5\) sur les douze points ne préserve l'ensemble des hexades de
\(W_{12}\) que pour l'élément identité. Ainsi,
l'entrelacement \(H_*\)-équivariant de \(W_{GK}\) ne se prolonge pas, avec ce
calibre, en une représentation commune sur \(E_{30}\).
\end{proposition}

\begin{proof}
Le premier énoncé est une séparation des domaines d'équivariance :
\(W_{GK}\) utilise \(H_*\subset A_5\), tandis que le théorème d'incidence
prouve l'équivariance de \(U_W\) pour
\(\operatorname{Aut}(W_{12})\). Le recensement exact reproductible applique les soixante
éléments de l'action \(A_5/C_5\) au système d'hexades étiqueté ; seule
l'identité conserve ce système. Comme \(H_*\) contient six éléments, son
action ne peut pas entrer par restriction de cette action fixe dans
\(\operatorname{Aut}(W_{12})\).
\end{proof}

La proposition précédente isole exactement le calibre fixe ; elle n'est pas une
obstruction abstraite au transport conjoint de la calibration. Le
propriétaire directeur de la recalibration \(S_3\)--Witt construit le monomorphisme et
la seconde flèche équivariants comme suit.

\begin{theorem}[Recalibration conjointe \texorpdfstring{\(H_*\)}{H*}--Witt]
\label{thm:recalibracion-hstar-witt}
Soit
\(\bar\rho_{\rm src}:H_*\to S_{12}\) l'action ponctuelle des douze classes
latérales déterminée par \(c_*\), \(r_*\) et l'isomorphisme \(\theta\) du
secteur orienté, et soit \(\rho_{\rm src}\) sa linéarisation dans \(V_{11}\).
Dans la convention des images indexées à partir de un, la permutation
\[
 \boxed{\sigma=(1,5,2,6,7,4,3,10,11,9,8,12)}
\]
définit le monomorphisme
\[
 \boxed{
 \iota_\sigma:H_*\hookrightarrow\operatorname{Aut}(W_{12})\cong M_{12},
 \qquad
 \iota_\sigma(h)=
 \sigma\,\bar\rho_{\rm src}(h)\,\sigma^{-1}.
 }
\]
Si la calibration est transportée intégralement au moyen de
\[
 W_{12}^{\sigma}=\{\sigma^{-1}B:B\in W_{12}\},
 \qquad
 \widetilde U_W:=U_WP_\sigma,
\]
où \(P_\sigma\) est l'opérateur unitaire de permutation, alors, sous
l'identification correspondante de l'espace d'incidence transporté à
\(E_{30}\), pour tout \(h\in H_*\),
\[
 \boxed{
 \widetilde U_W\,\rho_{\rm src}(h)
 =\rho_{E_{30}}\!\bigl(\iota_\sigma(h)\bigr)\,\widetilde U_W.
 }
\]
En particulier,
\[
 \widetilde Q_3=\widetilde U_WP_3\widetilde U_W^*
\]
est la réalisation de Witt du même projecteur et la chaîne
\[
 I_{L_{\rm or}}\otimes P_G
 \xleftrightarrow{\ W_{GK}\ }
 P_3
 \xleftrightarrow{\ \widetilde U_W\ }
 \widetilde Q_3
\]
porte, après identification de l'action source au moyen de \(\theta\), une
représentation commune de \(H_*\), avec l'action d'arrivée donnée par
\(\iota_\sigma\). La recalibration transporte conjointement
\(C_W,W_{12},I\), la dodécade, l'origine, l'orientation, la chambre et la
mémoire ; elle ne consiste pas à remplacer isolément \(U_W\).
\end{theorem}

\begin{proof}
Le certificat exact reconstruit le code de Golay ternaire, les \(132\)
hexades, le groupe \(M_{12}\) d'ordre \(95040\), le \(S_3\) source et un
\(S_3\) semi-régulier d'arrivée, tous deux d'ordre six. Il vérifie pour chaque
\(h\in H_*\) que
\(
\sigma\bar\rho_{\rm src}(h)\sigma^{-1}\in M_{12}
\), que \(W_{12}^{\sigma}\) est \(H_*\)-invariant et que son incidence
point--hexade est équivariante. L'injectivité de \(\iota_\sigma\) découle
du fait qu'il s'agit de la conjugaison de l'action fidèle source. De plus,
\[
 P_\sigma\rho_{\rm src}(h)
 =\rho_{12}\!\bigl(\iota_\sigma(h)\bigr)P_\sigma.
\]
En composant cette identité avec l'équivariance \(M_{12}\) de \(U_W\), on
obtient l'identité encadrée. Comme \(P_3\) commute avec \(H_*\), son
transport \(\widetilde Q_3\) commute avec \(\iota_\sigma(H_*)\) ;
l'équivariance de \(W_{GK}\) déjà démontrée complète la composition. Le même
certificat vérifie que le plan original n'est pas invariant sous l'action
source, de sorte que le no-go du calibre fixe est conservé sans contradiction.
\end{proof}

\textsc{Propriétaire matériel.} Ce résultat est régi par le théorème
\textsc{HMT--RDF--EXC--S3}, sections 3.1--3.3 du delta directeur de propagation
exceptionnelle Witt--Moonshine, et son certificat reproductible
\texttt{verificar\_recalibracion\_s3\_witt.py}, dont la sortie est
\texttt{PASS\_RECALIBRACION\_S3\_WITT}. Le statut est
\texttt{FORMALIZACION\_NUEVA/CERTIFICADO\_NUEVO} ; l'unicité du
changement de calibre n'est pas postulée.

Avant l'utilisation de \(B_K\), de \(u_K\), de l'orientation et du facteur polaire, il y avait
\(10\cdot6\cdot4=240\) applications étiquetées, ou \(40\) après identification des
réétiquetages internes du secteur. Les règles précédentes réduisent cet
ensemble à un élément. C'est le contenu vérifiable du mot
\emph{canonique} pour \(W_{GK}\), dans le cadre des règles déclarées. L'action
commune sur le module de Witt provient du changement explicite \(\sigma\) ; on
n'affirme pas que ce changement de calibre soit unique ou canonique.

La règle variationnelle \(\mathcal E_B\) n'utilise que des données HMT antérieures
et fixe la classe dans laquelle l'unicité est obtenue. Elle n'identifie pas l'involution
d'étoile à \(P_3\), car leurs spectres restent distincts.

\section{Comparaison des lectures de \texorpdfstring{\(t=7\) et \(t=8\)}{t=7 et t=8}}

En \(t=8\), il y a trois racines locales et onze extensions à l'horizon neuf de
l'unique racine qui se poursuit. Les tailles des trois fibres sont
\[
 (11,0,0).
\]
La matrice naturelle branche--chaîne a une ligne de onze uns et deux lignes
nulles, de sorte que son rang est un. En \(t=9\), de plus, le quotient visible
passe de \(3/11\) à \(3/6\). Cette fraction n'était pas la trace d'un
projecteur stable. La correction n'élimine pas le secteur : elle déplace la lecture
vers l'objet homogène \(P_G\) de \(t=7\), où le rang et la dimension
appartiennent au même espace.

\section{Conséquence pour une action positive}

La loi d'action peut maintenant utiliser une unique classe isométrique de projecteurs
positifs dans trois espaces de support :
\[
 P_G
 \xleftrightarrow{\ W_{GK}\ }
 P_3
 \xleftrightarrow{\ U_W\ }
 Q_3.
\]
L'étoile \(R_{B_K}\) reste réservée à la parité et à l'orientation. Le
projecteur \(P_G\) enregistre la survie ; \(P_3\) enregistre la polarisation
dodécaphasique ; \(Q_3\) enregistre sa réalisation d'incidence dans Witt. Cette
séparation des fonctions est la prémisse mathématique de la mesure et de la loi
positive de la section suivante. Dans le calibre fixe, la chaîne exprime
l'équivalence isométrique de projecteurs sans représentation commune. Le
\cref{thm:recalibracion-hstar-witt} transporte la calibration complète et
transforme la chaîne avec \(\widetilde U_W\) et \(\widetilde Q_3\) en une
représentation commune de \(H_*\), sans identifier l'étoile signée au
projecteur positif.
